% 102.literaturereview.tex

\section{Estado da Arte}
\label{sec:estado_arte}

\subsection{Trabalhos relacionados à detecção de nematoides}
\label{subsec:trabalhos_nematoides}

A detecção de nematoides fitoparasitas na soja emprega sensoriamento remoto,
visão computacional, análise química e sensores de gases. Essas abordagens
exploram respostas espectrais das plantas, imagens dos organismos e perfis
químicos associados às interações biológicas.

No sensoriamento remoto, \cite{wang2025scn} identificaram \textit{Heterodera
glycines} em plantas de soja por dados hiperespectrais e aprendizado de
máquina. O índice espectral SCNVI associado ao XGBoost obteve 70\% de acurácia
na classificação de três níveis de infestação. \cite{katari2026} avaliaram
redes neurais convolucionais (CNN) e Vision Transformers (ViT) em imagens
multiespectrais de áreas infestadas pelo nematoide de cisto da soja.
\cite{dellasilva2025} utilizaram dados hiperespectrais e classificadores para
avaliar níveis de infestação por \textit{Aphelenchoides besseyi} em soja. Esses
métodos dependem da resposta espectral da planta e não analisam gases
provenientes do solo.

Na análise por imagens microscópicas, \cite{park2026} avaliaram YOLO e
Detectron2 para detectar e segmentar fêmeas de \textit{H. glycines}. O YOLOv11
apresentou o melhor desempenho. A aplicação exige coleta das amostras, preparo
do material e aquisição das imagens. \cite{neupane2025} revisaram técnicas de
processamento de imagens e aprendizado de máquina para detecção e classificação
de nematoides. A revisão mostra o predomínio de métodos baseados em
características visuais.

Na análise química, \cite{rocha2023} investigaram solos associados à supressão
de \textit{H. glycines} em cultivos de soja e trigo. A extração por solvente e
a cromatografia gasosa acoplada à espectrometria de massas (GC--MS)
identificaram 51 compostos e diferenças entre os tratamentos. \cite{zhao2025}
avaliaram o efeito nematicida de \textit{Microbacterium maritypicum} Sneb159
contra \textit{H. glycines} e caracterizaram compostos voláteis bacterianos por
SPME-GC--MS. Esses estudos tratam de supressão e controle biológico, não da
detecção de solos infestados por sensores de gases.

No uso de nariz eletrônico, \cite{chen2026xtcn} combinaram sensores MOS, GC--MS
e uma rede convolucional temporal explicável, denominada X-TCN, para reconhecer
três espécies de insetos-praga de acervos culturais. O modelo obteve 99,28\% de
acurácia, 99,31\% de precisão, 99,26\% de recall e 99,27\% de F1. Os mapas de
atenção identificaram as regiões dos sinais com maior contribuição para a
classificação. O estudo demonstra o uso da X-TCN na análise temporal de sinais
de nariz eletrônico, mas não aborda nematoides, soja ou gases do solo.

\cite{herrmann2024} aplicaram sensores de gases e aprendizado de máquina à
identificação de estresse hídrico em plantas de soja. O modelo alcançou 94,4\%
de acurácia. \cite{kong2025} utilizaram nariz eletrônico, SVM e Random Forest
para detectar hidrocarbonetos de petróleo e pesticidas orgânicos em amostras
de solo. Esses estudos confirmam a aplicação de sensores de gases em problemas
agrícolas, mas não investigam nematoides fitoparasitas.

Os trabalhos mostram avanços na detecção de nematoides e no uso de narizes
eletrônicos. Métodos baseados em imagens dependem de características visuais ou
espectrais. Análises químicas exigem instrumentação laboratorial. Narizes
eletrônicos permitem a análise de perfis gasosos, mas os estudos encontrados
tratam de outros organismos, contaminantes ou condições de estresse.

Não foi localizado estudo que reúna nariz eletrônico, extração de gases sob
pressão forçada e variáveis ambientais para classificar solos cultivados com
soja com e sem \textit{H. glycines}. Esta pesquisa avalia essa combinação e
considera os efeitos da umidade do solo e da pressão nas respostas dos sensores
e no desempenho dos classificadores.

A Tabela~\ref{tab:trabalhos_relacionados_nematoides} compara as aplicações, as
tecnologias e as limitações dos estudos em relação a esta pesquisa.

% Impede que a tabela anteceda o texto da revisão.
\clearpage

\begin{table*}[!t]
\centering
\caption{Estudos relacionados à detecção de nematoides e ao uso de nariz eletrônico (2023--2026).}
\label{tab:trabalhos_relacionados_nematoides}

\renewcommand{\arraystretch}{1.20}
\setlength{\tabcolsep}{3.5pt}
\footnotesize

\begin{tabularx}{\textwidth}{@{}
>{\raggedright\arraybackslash}p{2.15cm}
>{\centering\arraybackslash}p{0.85cm}
>{\raggedright\arraybackslash}p{2.45cm}
>{\raggedright\arraybackslash}p{2.85cm}
>{\centering\arraybackslash}p{0.95cm}
>{\raggedright\arraybackslash}X
@{}}

\toprule
\textbf{Autor} &
\textbf{Ano} &
\textbf{Aplicação} &
\textbf{Tecnologia} &
\textbf{ML} &
\textbf{Limitação frente à proposta} \\
\midrule

Chen et al. \cite{chen2026xtcn}
& 2026
& Insetos-praga de acervos culturais
& Sensores MOS, GC--MS, X-TCN e IA explicável
& Sim
& Não avalia nematoides, soja ou gases do solo \\

Katari et al. \cite{katari2026}
& 2026
& \textit{H. glycines} em soja
& Imagens multiespectrais, CNN e ViT
& Sim
& Não analisa gases provenientes do solo \\

Park et al. \cite{park2026}
& 2026
& \textit{H. glycines} em solo
& Microscopia RGB, YOLO e Detectron2
& Sim
& Exige preparo das amostras e microscopia \\

Wang et al. \cite{wang2025scn}
& 2025
& \textit{H. glycines} em soja
& Hiperespectral, SCNVI e XGBoost
& Sim
& Depende da resposta espectral da planta \\

Della-Silva et al. \cite{dellasilva2025}
& 2025
& \textit{A. besseyi} em soja
& Hiperespectral e classificadores
& Sim
& Outra espécie; não analisa gases do solo \\

Rocha et al. \cite{rocha2023}
& 2023
& Solos supressivos a \textit{H. glycines}
& Extração por solvente e GC--MS
& Não
& Caracterização química, sem diagnóstico por sensores \\

Zhao et al. \cite{zhao2025}
& 2025
& VOCs bacterianos contra \textit{H. glycines}
& SPME-GC--MS e bioensaios
& Não
& Avalia atividade nematicida, não detecção \\

Herrmann et al. \cite{herrmann2024}
& 2024
& Estresse hídrico em soja
& Nariz eletrônico e ML
& Sim
& Não avalia infestação por nematoides \\

Kong et al. \cite{kong2025}
& 2025
& Contaminantes químicos no solo
& Nariz eletrônico, SVM e RF
& Sim
& Não avalia patógenos do solo \\

Neupane et al. \cite{neupane2025}
& 2025
& Detecção de nematoides (revisão)
& Visão computacional e ML
& Rev.
& Ênfase em dados de imagem \\

\bottomrule

\end{tabularx}
\end{table*}
